\documentclass{article}

% ICML 2022 style approximation
\usepackage[letterpaper, top=1in, bottom=1in, left=1.25in, right=1.25in]{geometry}
\usepackage{times}
\usepackage{microtype}
\usepackage{graphicx}
\usepackage{hyperref}
\usepackage{url}
\usepackage{booktabs}
\usepackage{amsmath}
\usepackage{amssymb}
\usepackage[numbers]{natbib}
\usepackage{titlesec}
\usepackage{parskip}
\usepackage{authblk}
\usepackage{abstract}

% ICML-style section formatting
\titleformat{\section}{\normalfont\large\bfseries}{\thesection.}{0.5em}{}
\titleformat{\subsection}{\normalfont\normalsize\bfseries}{\thesubsection.}{0.5em}{}
\titlespacing*{\section}{0pt}{1.2ex plus 1ex minus .2ex}{0.8ex plus .2ex}
\titlespacing*{\subsection}{0pt}{0.8ex plus 1ex minus .2ex}{0.5ex plus .2ex}

% Compact spacing
\setlength{\parskip}{0.5ex}
\setlength{\parindent}{0pt}

% Abstract styling
\renewenvironment{abstract}
 {\small\begin{center}\bfseries Abstract\end{center}\begin{quotation}}
 {\end{quotation}}

\begin{document}

% Title
\begin{center}
{\LARGE \textbf{Do Modern Deep Learning Architectures Disproportionately\\
Benefit BCI-Inefficient Users? An Empirical Analysis\\
on Motor Imagery EEG Classification}}

\vspace{1em}

{\large Chris Dollo \quad Madeline Rippin}\\[0.3em]
{\normalsize CMSC 678 --- Introduction to Machine Learning\\
University of Maryland, Baltimore County}\\[0.3em]
{\small \today}
\end{center}

\vspace{1em}

\begin{abstract}
Brain-computer interfaces (BCIs) based on motor imagery electroencephalography (EEG) offer transformative potential for assistive technologies, yet 15--30\% of users cannot effectively operate them --- a phenomenon termed BCI inefficiency. Prior work has shown that convolutional neural networks (CNNs) improve classification for inefficient users more than traditional signal processing methods, but this has been studied only for binary motor imagery tasks using online BCI accuracy as the performance criterion. We extend this investigation to 4-class motor imagery using the BCI Competition IV Dataset 2a, comparing three models spanning the pipeline from traditional to state-of-the-art: Filter Bank Common Spatial Pattern (FBCSP), EEGNet, and a Mamba-based state space model. We define high and low performer groups via median split on FBCSP kappa, then analyze whether gains over the traditional baseline are disproportionately larger for low performers. We further investigate whether simple data augmentation strategies (Gaussian noise addition, sliding window) differentially benefit low-performing subjects. Our work provides the first systematic analysis of BCI inefficiency across this model hierarchy in 4-class classification, with direct implications for the equitable deployment of BCI-based assistive and rehabilitative technologies.
\end{abstract}

\section{Problem Description and Motivation}

Motor imagery EEG classification is a core task in brain-computer interface (BCI) research, with critical applications in neuroprosthetics, post-stroke rehabilitation, and assistive communication for individuals with motor disabilities. A persistent and clinically significant obstacle is BCI inefficiency: an estimated 15--30\% of users are unable to produce the desired sensorimotor rhythm (SMR) patterns robustly enough for a classifier to exceed 70\% accuracy, even after extensive training~\cite{blankertz2009bci,ahn2013theta}. This is not merely an accuracy problem --- it is an equity problem. The patients who most urgently need BCI technology (e.g., ALS patients, stroke survivors, upper-limb amputees) are often precisely those who are harder to classify, whether due to neurological disruption, limited motor cortex plasticity, or simply natural inter-individual variability in EEG signal quality. Recent deep learning models have dramatically improved average accuracy on benchmark datasets like BCI Competition IV Dataset 2a~\cite{tangermann2012review}, but nearly all published work reports mean accuracy across subjects without examining whether these gains are distributed equitably. Whether modern architectures specifically help the users who need BCIs most remains an open and practically important question.

\section{Research Questions}

We intend to examine the following research questions:

\textbf{RQ1:} Does EEGNet improve classification accuracy for low-performing (BCI-inefficient) subjects more than for high-performing subjects, relative to the FBCSP baseline? This extends the binary-task finding of~\cite{nakanishi2022classification} to 4-class motor imagery on BCI IV-2a.

\textbf{RQ2:} Does a Mamba-based state space model (SSM) further amplify this effect relative to EEGNet? The theoretical motivation is that Mamba's linear-complexity long-sequence modeling may be better suited to capturing weak, diffuse, or temporally irregular ERD/ERS patterns that characterize low-performing subjects.

\textbf{RQ3:} Does data augmentation (Gaussian noise injection and sliding window) differentially benefit low-performing subjects across all three models? Specifically, we ask whether augmentation closes the performance gap between high and low performers, or merely shifts the mean.

Together, these questions address both the architectural dimension (which model family helps hard subjects most?) and the data dimension (does more training data help hard subjects more than easy ones?), producing a comprehensive empirical picture of BCI inefficiency across the model hierarchy.

\section{Proposed Solution}

We implement and compare three models on BCI Competition IV Dataset 2a. Our \textbf{traditional baseline} is FBCSP with LDA, using the one-vs-rest multiclass extension~\cite{ang2012filter}. This replicates the published competition-winning approach. Our \textbf{CNN baseline} is EEGNet~\cite{lawhern2018eegnet} configured per the original paper (128~Hz resampling, temporal kernel length~64, session 1 train/session 2 test protocol). Our \textbf{novel model} is a Mamba-based SSM, specifically MI-Mamba~\cite{guo2025mimamba}, a hybrid CNN-Mamba architecture that processes local spatial features via convolution before applying Mamba's selective state-space mechanism for global temporal modeling. All three models are evaluated under identical preprocessing and train/test protocols to ensure fair comparison.

Subject performance groups are defined via a median split of per-subject FBCSP kappa across the 9 subjects~\cite{nakanishi2022classification}. This yields five high performers and four low performers, consistent with known per-subject rankings in the literature. The primary analysis computes $\Delta\text{Acc} = \text{Acc}_\text{DL} - \text{Acc}_\text{FBCSP}$ for each subject, converted to a common accuracy scale, and compares the mean $\Delta\text{Acc}$ between groups using a Mann-Whitney U test (appropriate for the small $n=4$ vs $n=5$ comparison). We additionally perform spectral analysis of low vs.\ high performers --- computing power spectral density and topographic maps to characterize the neurophysiological profile of hard subjects in terms of mu/beta ERD strength.

For the augmentation experiment, we apply two strategies to session 1 training data only: (1) \textbf{Gaussian noise injection} producing 3 augmented copies per trial by adding zero-mean Gaussian noise scaled to 5--10\% of each trial's standard deviation, yielding $288 \times 3 + 288 = 1152$ training trials per subject; and (2) \textbf{sliding window segmentation} with a 2-second window and 50\% overlap, generating approximately 3--4 windows per 4-second trial. We train all three models under both the original and augmented conditions and repeat the BCI illiteracy group analysis, asking whether augmentation differentially improves accuracy for low performers.

\section{Relation to Prior Work}

The most directly related work is~\cite{nakanishi2022classification}, which found that a CNN trained on raw EEG significantly outperformed CSP+LDA for inefficient users in a 2-class online BCI paradigm, with improvements of up to 28\% for the lowest performers. Our work extends this to 4-class offline classification on BCI IV-2a and introduces the Mamba architecture as a third point of comparison. The FBCSP baseline on BCI IV-2a is established by~\cite{ang2012filter}, who reported per-subject kappa values for all 9 subjects; EEGNet on this dataset has been evaluated by multiple groups~\cite{lawhern2018eegnet, altaheri2023atcnet}. The BCI illiteracy phenomenon is documented in~\cite{blankertz2009bci} and~\cite{ahn2013theta}, the latter showing that high theta and low alpha power in resting EEG may predict poor BCI performance. Mamba-based models for MI-EEG have appeared very recently~\cite{guo2025mimamba, yang2024stmamba}, reporting within-subject accuracy of 80--87\% on BCI IV-2a, but none analyze per-subject performance through the BCI inefficiency lens.

Data augmentation for EEG motor imagery has been studied in~\cite{zhao2023augmentation}, which found that Gaussian noise and sliding window augmentation improve accuracy on BCI IV-2a, and~\cite{xie2024bfatcnet}, which confirmed that Gaussian noise augmentation improves noise immunity. Neither work examines differential effects on BCI-inefficient subjects specifically.

\section{Anticipated Challenges}

The primary practical challenge is the small sample size inherent to BCI IV-2a: only 9 subjects, yielding groups of 4 and 5 after the median split. This severely limits statistical power, and the Mann-Whitney U test at $n=4$ vs $n=5$ has limited ability to detect small-to-moderate effect sizes. We will be explicit about this limitation in our final report and interpret results cautiously, framing them as exploratory rather than confirmatory. Implementing MI-Mamba requires a working Mamba installation, which depends on CUDA-specific kernels that may not be available in all compute environments; if this proves intractable, we will fall back to STMambaNet, which has a more portable implementation.

A conceptual challenge is that BCI IV-2a was not collected with online BCI accuracy scores, unlike the~\cite{nakanishi2022classification} study, so our group assignment relies on FBCSP kappa as a proxy rather than true BCI literacy. We will validate this proxy by cross-checking against known per-subject rankings in the literature and discussing its limitations explicitly. Finally, ensuring that data augmentation is applied correctly --- only to training data and not contaminating the test set --- requires careful pipeline management that we will document thoroughly.

\section{Ethical Components}

This project directly engages with a significant equity concern in assistive technology: BCI systems perform worst for the users who need them most. \cite{blankertz2009bci} reframed ``BCI illiteracy'' as a systems failure rather than a user failure, and our work operationalizes this view empirically. If certain architectures or augmentation strategies disproportionately help low performers, this has direct implications for model selection in clinical deployment --- a finding with real stakes for patients with ALS, spinal cord injury, or post-stroke motor impairment who may have limited access to lengthy calibration sessions.

The dataset used (BCI IV-2a) was collected from healthy volunteers under controlled research conditions, so direct participant harm is not a concern. However, any model deployed in a clinical BCI application based on this research would require validation on clinical populations, which have different EEG characteristics than healthy subjects and are not represented in this dataset. We will discuss this generalization gap explicitly. We will also note that algorithmic improvements in BCI classification, while beneficial, do not replace the need for broader accessibility design in BCI hardware and software systems.

\bibliography{references}
\bibliographystyle{plain}

\end{document}
